\section{Conclusion}
\label{sec:sonuc}

This study systematically examined the shape of the blending function
used where a coarse-resolution global terrain meets a fine-detail local
patch in an Isaac Lab-based lunar-landing environment, both at the
analytic/CPU level and in a real Isaac Sim SAC training run. Analytic
measurements over thirty seeds statistically confirmed that the
previously used as-shipped Gaussian blend produces, once truncated at a
finite radius, boundary continuity worse than smoothstep by three
orders of magnitude, and that this defect is a scale-invariant,
structural property. The proposed hybrid method and a calibrated
Gaussian variant remove this defect to near-machine precision --- this
is the study's most robust and least debatable finding.

The real Isaac Sim training comparison began with a single seed per
method, and the initial measurement showed Gaussian and hybrid reaching
a markedly higher success rate than smoothstep. To test the reliability
of this finding, the study was extended with two additional independent
seeds (total $n=3$), and no statistically significant difference was
found between the methods on any metric. This is the study's second
major contribution: it provides a concrete, self-contained example of
how misleading single-run RL comparisons can be, and empirically
confirms the warning of Henderson et al.\ \citep{henderson2018deep} in
this specific setting. Taken together with the analytic results, this
shows that the choice of blending function has a measurable and robust
effect on boundary continuity, but that this effect --- at least at the
present scale and with $n=3$ --- is not reflected in final RL training
performance in a distinguishable way.

The core contribution of this study is a methodology that tests the
choice of blending function quantitatively, both analytically and at the
level of RL training, in a real Isaac Sim/Isaac Lab environment, rather
than treating it as a purely visual/rendering preference; this
methodology is able to reliably report a ``no significant difference''
result just as well as a positive one. All three methods are integrated
into the production code as a backward-compatible, selectable
configuration, making these findings directly actionable: with the
current data, there is no statistical justification for choosing among
the three methods on RL training performance; the choice can instead be
based on other criteria such as boundary continuity (analytically in
favor of Gaussian and hybrid) or computational cost (in favor of
smoothstep).

\subsection*{Implications for autonomous control system design}
For engineers designing the training pipeline of an autonomous control
system, the key takeaway is not which blending function to prefer but
that a simulator's internal geometry-generation choices are themselves
part of the control system's design space and warrant the same
quantitative scrutiny as the controller or the reward function. A
world-model artifact with a large, statistically confirmed effect on
boundary continuity (here, close to three orders of magnitude) need not
automatically threaten the learned controller's deployed behavior; this
study shows a case where a robust simulation-fidelity defect and a
distinguishable control-performance regression can, at a given training
scale, come apart. The practical implication is that such simulator
artifacts should be audited and corrected on their own engineering
merits --- because a discontinuity in a physically queryable contact
geometry is exactly the kind of defect a real vehicle's contact sensors
would not tolerate --- rather than assumed safe merely because a single
training run did not reveal a resulting control-performance gap.

Three concrete directions are proposed for future work. First, because
$n=3$ still has low statistical power, the real-training comparison
should be extended with at least five to ten more independent seeds to
determine conclusively whether the small observed average differences
(particularly the Gaussian's consistently highest average) reflect a
real effect or residual noise. Second, the training scale should be
extended to the full configuration (1200 iterations) to confirm whether
the observed plateaus reflect an interim-stage rather than a
final/asymptotic performance difference. Third, going beyond the scope
that this study deliberately kept narrow (see \cref{sec:tartisma}), the
external validity of the findings should be tested: repeating the same
three blending methods on a different physics engine, a different
reward/observation design, or a different contact-rich task (e.g.,
quadrupedal locomotion, humanoid standing) would disentangle whether the
findings here depend on the blending function or on other design choices
specific to this particular environment.
